\Askhsh{28342}{4}
\begin{ylh}{1}
    \item 1.8 Συνέχεια συνάρτησης
    \item 2.6 Συνέπειες του Θεωρήματος Μέσης Τιμής
    \item 2.7 Τοπικά ακρότατα συνάρτησης
    \item 2.9 Ασύμπτωτες - Κανόνες De l’ Hospital.
\end{ylh}
\begin{wrapfigure}[20]{r}{7.5cm} % The [20] is the number of lines of text to wrap around. Adjust as needed.
\centering
\begin{tikzpicture}
\begin{axis}[
    axis lines=middle,
    xlabel=$x$,
    ylabel=$y$,
    xmin=-0.5, xmax=2.7, % Adjusted xmin/xmax for better fit
    ymin=-0.2, ymax=3, % Adjusted ymin/ymax
    xtick=\empty, % Custom labels for x-axis points
    ytick=\empty, % No y-ticks needed explicitly from image
    grid=none, % No grid in the image
    samples=100,
    width=7cm, height=6cm,
    tick label style={font=\scriptsize},
    every axis x label/.style={at={(current axis.right of origin)}, anchor=north west},
    every axis y label/.style={at={(current axis.above origin)}, anchor=south},
    axis line style={-stealth}, % Arrows on axes
]
    % Function y = e^x for x < 1
    \addplot[very thick, black, domain=-0.5:1] {exp(x)};
    \node[black, below left] at (axis cs:0.2, 1.2) {$y = e^x$};

    % Function y = e/x for x > 1
    \addplot[very thick, black, domain=1:2.7] {exp(1)/x};
    \node[black, above right] at (axis cs:1.8, 1.5) {$y = \frac{e}{x}$};

    % Dashed line at x=1
    \draw[dashed, black] (axis cs:1,0) -- (axis cs:1,exp(1));
    \node[circle,draw,inner sep=1pt,fill=white, black] at (axis cs:1,exp(1)) {}; % Hollow circle at (1,e)

    % Rectangle ABΓΔ
    % Use α = 0.4 for drawing as in the image (O is slightly left of A)
    \pgfmathsetmacro{\alpha_val}{0.4}
    \pgfmathsetmacro{\xdelta_val}{exp(1-\alpha_val)} % x_Δ = e^(1-α)
    \pgfmathsetmacro{\y_val}{exp(\alpha_val)} % Height = e^α

    % Coordinates for the rectangle
    \coordinate (A_pt) at (axis cs:\alpha_val, 0);
    \coordinate (B_pt) at (axis cs:\alpha_val, \y_val);
    \coordinate (Gamma_pt) at (axis cs:\xdelta_val, \y_val);
    \coordinate (Delta_pt) at (axis cs:\xdelta_val, 0);

    % Draw the rectangle
    \draw[very thick, fill=gray!30] (A_pt) -- (B_pt) -- (Gamma_pt) -- (Delta_pt) -- cycle;

    % Points and labels on the x-axis
    \fill (axis cs:0,0) circle (1.5pt); % O
    \node[below] at (axis cs:0,0) {O};

    \fill (A_pt) circle (2pt); % A
    \node[below right] at (A_pt) {A($\alpha$,0)}; % Adjust position to avoid overlap with O

    \draw (axis cs:1,0) circle (2pt); % 1 (hollow circle)
    \node[below] at (axis cs:1,0) {1};

    \fill (Delta_pt) circle (2pt); % Δ
    \node[below] at (Delta_pt) {$\Delta$};
    
    % Points and labels for B and Γ
    \fill (B_pt) circle (2pt);
    \node[left] at (B_pt) {B};
    
    \fill (Gamma_pt) circle (2pt);
    \node[right] at (Gamma_pt) {$\Gamma$};
    
    % x' label on the left of the x-axis
    \node[below] at (axis cs:-0.5,0) {$x'$};
\end{axis}
\end{tikzpicture}
\end{wrapfigure}
Στο παρακάτω σχήμα το ορθογώνιο ΑΒΓΔ έχει τις κορυφές Α και Δ πάνω στον άξονα $x'x$ και τις κορυφές Β και Γ πάνω στις γραφικές παραστάσεις των συναρτήσεων $f(x) = e^x \, , \, x < 1$ και $g(x) = \frac{e}{x} \, , \, x > 1$, αντίστοιχα. Έστω A$(\alpha, 0)$ με $\alpha < 1$.

\begin{alist}
\item Να αποδείξετε ότι:
    \begin{rlist}
    \item η τετμημένη της κορυφής $\Delta$ είναι $x_\Delta = e^{1-\alpha}$, \monades{6}
    \item το εμβαδόν του ορθογωνίου ΑΒΓΔ είναι $E(\alpha) = e - \alpha e^\alpha$, $\alpha < 1$. \monades{6}
    \end{rlist}
\item Να βρείτε τη μέγιστη τιμή του εμβαδού του ορθογωνίου ΑΒΓΔ. \monades{7}
\item Να εξετάσετε αν υπάρχουν και πόσες τιμές του $\alpha$, για τις οποίες το εμβαδόν του ορθογωνίου ΑΒΓΔ γίνεται ίσο με 1. \monades{6}
\end{alist}